\paragraph{The residual positive complex.}
Write $K_n=\mathbf Q(\zeta_{2^{n+2}})^+$, and let $S_n$ consist of
the places above two and a fixed finite set of odd rational primes.
Put $r_n=2^n$ and $s_n=\#S_n$. We first compute with the scalar
coefficient $\mathbf F_2=\mu_2$.

The ordinary two-class group of $K_n$ is trivial. Indeed this holds
for $K_0=\mathbf Q$. Assuming it for $K_{n-1}$, let $H$ be the
Hilbert two-class field of $K_n$, with the real places split.
The canonical extension $H/K_{n-1}$ is Galois with a finite
two-group $G_H$. Only the unique dyadic prime can ramify, and its
inertia $I$ has order two. The quotient of $G_H^{\rm ab}$ by the image
of $I$ gives an everywhere unramified abelian two-extension of
$K_{n-1}$ with real places split, hence is trivial by induction.
Thus $G_H^{\rm ab}$ is cyclic, generated by an element of order at most
two. The Frattini quotient of $G_H$ has one generator, so $G_H$ is cyclic and
has order at most two. Since it contains $K_n$, we obtain $H=K_n$.
The $S_n$-class group is a quotient, so its two-part also vanishes.

The signatures of units span all $r_n$ real signs. For $n\ge1$ put
$m=n+2$, $\zeta=\zeta_{2^m}$, and
\[
 u=1+\zeta+\zeta^{-1}
   =\zeta^{-1}\frac{\zeta^3-1}{\zeta-1}.
\]
This is a unit: the numerator and denominator generate the same
dyadic prime ideal in the cyclotomic field. Its real conjugates are
$1+2\cos(2\pi a/2^m)$ for odd $0<a<2^{m-1}$.
The negative ones have $2^m/3<a<2^{m-1}$, and their number is
\[
 2^{m-2}-\left\lfloor\frac{2^{m-1}}3+\frac12\right\rfloor,
\]
an odd integer. Hence $N_{K_n/\mathbf Q}(u)=-1$.
The signature vector of $u$ has augmentation one in the local ring
$\mathbf F_2[\operatorname{Gal}(K_n/\mathbf Q)]$ and is therefore
a unit. Its Galois translates span all signatures. At $n=0$ use $-1$.

Kummer theory now gives
$\dim H^1(G_{K_n,S_n},\mathbf F_2)=r_n+s_n$, and its signature map
to $\mathbf F_2^{r_n}$ is onto. The map in degree zero is the diagonal
injection. The positive complex is the fiber of ordinary global
cochains to the sum of ordinary real cochains, so its long exact
sequence gives $\dim H^1(C_n^+)=r_n+s_n-1$.
The vanishing two-part of the $S_n$-class group identifies global
$H^2$ with $\operatorname{Br}(\mathcal O_{K_n,S_n})[2]$.
The local invariant description has $r_n+s_n$ coordinates with one
sum relation. Projection to the $r_n$ real coordinates is onto,
because $s_n\ge1$, and has kernel of dimension $s_n-1$.
Together with signature surjectivity this gives
\[
 \dim H^1(C_n^+)=2^n+s_n-1,\qquad
 \dim H^2(C_n^+)=s_n-1,
 \qquad H^i(C_n^+)=0\quad(i\ne1,2).
\]
The last assertion uses the usual global-to-real comparison in
degrees greater than two; degree zero vanishes by the diagonal injection.

